\subsection{Transporting the pointwise bounds at the prime assigned to the sequence}
\label{subsec:word-period}

The preceding orbit estimates are away from the prime assigned to the sequence. The coefficient places above that prime fall into two classes, with different support conditions.

Let \(T=T_u=\pi\bar\pi\) be the chosen prime assigned to the sequence and let \(w\mid T\)
have ramification index \(e_w\). Call the place distinguished if it
divides the embedded \(\pi\). Write \(V_w\) for the local factors
corresponding to tops under the position bijection at a distinguished
place, and \(V_w=\varnothing\) at the conjugate places. Set
\(s=(t+1)/2\), 
\[
  N_i=k+i(t+1).
\]
 Let \(S_i\) be the integral space
of values of the twisted forms of level \(i\) at \(w\), and \(K_i\)
the transition to level \(i+1\).

For this subsection its pointwise
kernel is denoted by \(k_i(P,P')\). For a set \(V\) of factors
define \(B_i(P,V)\) by the integral coefficient tests in mixed
frames. At a distinguished place the frame uses the kernel of the
residue row at factors in \(V\) and the underlying lattice elsewhere.
At a conjugate place it uses the kernel frames with the conjugate
coefficient action at every factor. On \(V=\varnothing\) the tests
at the representatives define \(S_i\). All these tests are
independent of the choice of adapted bases by their integral
unimodular transition matrices.

\begin{proposition}[Valuation after transporting the decorated column]\label{prop:word-prime-transfer}
Suppose the pointwise estimates \eqref{eq:tag-w1} and incidence identity
\eqref{eq:tag-w2} hold at every pair of points at the assigned levels.
On the actual local trace-zero column, the composite transition
to the terminal lattice has denominator killed by every
coefficient \(c\) with \(v_w(c)\ge e_w m E_T(w)\).
Here \(E_T(w)=0\) at a conjugate place. At a distinguished
place it is
\[
 E_T(w)=2\sum_\tau W_\tau+fN_m+
              |V_w|\frac{(t+1)(m+1)}4.
\]
This applies to each twisting tuple satisfying the same
pointwise hypotheses.
\end{proposition}

\begin{proof}
\proofstep{Set out the two local inclusions and the neighbor incidence identity.}
Here are the exact point estimates on which the estimate at the prime assigned to the sequence
depends. They are needed at every pair of points at the indicated
levels. Denote the complete set of factors by \(V_{\rm all}\).
At a distinguished place and a conjugate place, respectively, they read
\[
\begin{split}
 T^{fN_i+\sum_\tau s_{{\rm div},i}(\tau)}
    k_i(P,P')B_i(P,\varnothing)
       &\subset B_{i+1}(P',V_{\rm all}),\\
    k_i(P,P')B_i(P,\varnothing)
       &\subset T^{fs} B_{i+1}(P',\varnothing).
\end{split}                                                       \tag{W1}\label{eq:tag-w1}
\]
The pointwise correspondence for the decorated neighbors at \(v\) is
\[
  \sum_{P'_1:\,P'\to_v P'_1}k_i(P,P'_1)y
       =T^s\sum_{P_0:\,P_0\to_v P}k_i(P_0,P')y .               \tag{W2}\label{eq:tag-w2}
\]
The two sums in \eqref{eq:tag-w2} run through individual neighbors before passage
to classes. We prove \eqref{eq:tag-w2} directly. Use \((D,\ell)\) and
\((D',\ell')\) for the two local lattices with surjective residue
rows, and take \(z_0\in D\) with \(\ell(z_0)=1\). The source
marking has kernel \(H\). The finite reader in this factor takes
value one on \(x\) precisely when \(xD\subset D'\) and
\(\ell'x=\ell\), and zero otherwise.

An upper target neighbor corresponds to \(y\in D'/TD'\) with
\(\ell'(y)=1\), with lattice \(D'_y=D'+\mathcal O y/T\) and
row \(z\mapsto\ell'(Tz)\). The latter row has kernel \(D'\).
A lower source neighbor has lattice \(H\) and row \(h\) on
\(H/TH\) with \(h(Tz_0)=1\).

For a rational \(x\), either
sum of indicators is one precisely when
\[
                xH\subset D',\qquad Txz_0\in D',
                  \qquad \ell'(Txz_0)=1.
\]
Indeed at an admissible upper neighbor, the row equality on
\(H\) forces \(xH\subset D'\), and the row on \(z_0\) forces
\(y=Txz_0\bmod TD'\). Conversely this \(y\) works because
\(D=H+\mathcal O z_0\) and \(\ell\) is the coefficient of
\(z_0\) modulo \(T\). At a lower neighbor the only possible
row is \(h=\ell'x|_H\), and it works under exactly the three
displayed conditions. Thus the test includes every map, of
any rank. All other local decorations are fixed, since the
prime assigned to the sequence is disjoint from the selected twisting primes.

For each upper summand move both underlying companions down
by one local index. They become the companions of \(H,D'\),
which are the same for every lower summand. The split vacuum
comparison on a fixed away reader at source and target
determinant increments \(d,d'\) has scalar
\[
          \omega_v^{d'-d} T^{N_{i+1}d/2-N_i d'/2}.
\]
Here the unramified step and twisting characters give the
value \(\omega_v\) at the uniformizer. The other factor is
the inverse square root of the ratio of the two Hom volumes:
for a source index one the Hom ball expands by
\(T^{N_{i+1}}\), and for a target index one it contracts
by \(T^{N_i}\).

To make the comparison for this reader,
pass to a common refinement of its two finite translation
quotients. For a coset \(C\) partitioned into its finer cosets \(C'\), the identity of theta sections follows term by term from
\[
  \mathbf1_C=\sum_{C'\subset C}\mathbf1_{C'}.
\]
 The away values and derivatives agree there,
and the two split lattice changes act on the same indicator.
The reader on the refined quotient is still the product of
the preceding local test and its fixed away part. Thus the
move has \(d=d'=1\) at this factor and zero at the others,
so its scalar is \(T^s\). Linearity and the indicator equality
prove \eqref{eq:tag-w2}. Target equivariance follows by changing the
target variable in the theta sum and in its frame polynomial.
\label{cor:word-period-1}

Extend integral values from the level representatives by

\[
  \varphi(\gamma P_b)=\gamma\varphi(P_b).
\]
 The stabilizer
condition makes this single valued, and transport preserves
the adapted integral frames. In the kernel sum each denominator
\(|\operatorname{Stab}P_b|\) is a \(T\)-unit. Otherwise Cauchy's
theorem would give an element \(\gamma\) of order \(T\), and the
faithful rational representation would imply
\[
 T-1=\deg\Phi_T\le \deg\min_{\mathbb Q}(\gamma)
                         \le 2fN_m<T-1 .
\]
The sum therefore lies in the target lattice by \eqref{eq:tag-w1}.
Target equivariance extends that conclusion from representatives
to every point.
For an integral functional \(f\) and a unit-valued level character \(\chi\), a change of level representative preserves the test because
\[
  v_w(\chi(h)f)=v_w(f) .
\]
At a distinguished place the operator \(U_{i+1,v}\) removes \(v\)
from the set of mixed factors. For each forward neighbor at the
target, the kernel frame of that neighbor at \(v\) spans the
ordinary lattice of the original target. The other frames have
the original spans. Thus every summand passes the reduced test.

At an unused factor \(v\) one can remove the mixed condition
by multiplying instead by \(T^{\lambda_1(\tau(v))}\), where
\(\lambda\) is the true target weight. In adapted bases the change of frame is diagonal up to integral invertible factors, with elementary divisors
\[
  T,\,1,\ldots,1.
\]
Thus the least elementary divisor of its inverse has valuation \(-1\), in the valuation giving \(T\) order one. For a
polynomial representation with bottom weight removed the
maximum degree in the changed coordinate is the number of
columns of the diagram, since each minor uses it at most once.
Restoring the determinant power adds the bottom weight to this
degree and gives \(\lambda_1\).

If the bottom weight is \(\lambda_k\), the polynomial width and the determinant power give
\[
  (\lambda_1-\lambda_k)+\lambda_k=\lambda_1 .
\]
 The argument also applies to
negative bottom weights. Consequently multiplying by this power
puts the mixed value in the smaller test with \(v\) removed.
Use the \(U\) operation at the tops and the scalar operation at
every other factor. The sum of the latter exponents is at most
\(\sum_\tau W_\tau\), since each first weight off the tops is
bounded by its \(W_\tau\). Also
\(\sum_\tau s_{{\rm div},i}(\tau)\le\sum_\tau W_\tau\).
Let \(U_j\) be the ordered product at \(V_w\). We get
\[
 T^{e} U_{i+1}K_i S_i\subset S_{i+1},\qquad
 e=\begin{cases}
       2\sum_\tau W_\tau+fN_m,&w\mid\pi,\\
       0,&w\nmid\pi ,
     \end{cases}                                                  \tag{W3}\label{eq:tag-w3}
\]
In the second
case \eqref{eq:tag-w1} already gives ordinary integral values and the empty
product acts as the identity.

To transfer the \(U\)'s across \(K_i\), apply \eqref{eq:tag-w2} term by term
to the equivariant extension. Finite incoming and outgoing
fibers give two ways of summing over orbits of incident pairs
\((P_0,P)\). For a fixed marked first endpoint its stabilizer
acts on the finite fiber, and the reciprocal order of that
stabilizer times the orbit size is
\(1/|\operatorname{Stab}(P_0,P)|\). Fixing the second endpoint
gives exactly the same mass. This double count and source
covariance yield
\[
 U_{i+1,v} K_i=T^s K_i U_{i,v},\qquad
 U_{i+1}K_i=T^{s|V_w|}K_i U_i .                                 \tag{W4}\label{eq:tag-w4}
\]

All the sums act on forms: the neighbor sets are equivariant,
and the point kernel is equivariant at the target. Induction
moves one full product \(U_j\) through \(j\) successive
transitions with scalar \(T^{js|V_w|}\). Applying \eqref{eq:tag-w3} at each
stage therefore proves
\[
 T^{b}(K_{m-1}\cdots K_0)U_0^m S_0\subset S_m,\qquad
 b=me+s|V_w|\frac{m(m+1)}2=m E_T(w).                              \tag{W5}\label{eq:tag-w5}
\]
Equivalently, \(E_T(w)=0\) at a conjugate place and at a
distinguished place
\[
 E_T(w)=2\sum_\tau W_\tau+fN_m+
             |V_w|\frac{(t+1)(m+1)}{4}.
\]
In particular one pays the commutation at stage \(j\) exactly
\(j\) times. The weaker bound obtained by charging \(m\) times
at every stage need not meet the specified price.

\proofstep{Use the factor produced by repeated commutation on the actual column.}
The actual local column is the image of \(S_0\cap R\) under
\(T^{a}U_0^m\), with \(a=m\sum_{v\in V_w} b_v\ge0\) as in
Section~\ref{sec:ramified-column}. Transporting the initial
untwisted space through its level-zero identification preserves
the mixed tests, top weights and neighbor operators. For every
coefficient \(c\) with \(v_w(c)\ge e_w b\), the element
\(cT^{a-b}\) is integral at \(w\). Multiplying \eqref{eq:tag-w5} by this
element proves that \(c\) times the true column maps into the
terminal lattice. This gives the contribution at the assigned prime for
every twisting tuple satisfying the same point bounds.
\label{cor:word-period-2}
\end{proof}

